\documentclass[11pt,a4paper]{article}

% -------------------------------------------------
% Sprache / Encoding
% -------------------------------------------------
\usepackage[T1]{fontenc}
\usepackage[utf8]{inputenc}
\usepackage[ngerman]{babel}
\usepackage{lmodern}
\usepackage{microtype}

% -------------------------------------------------
% Mathematik
% -------------------------------------------------
\usepackage{amsmath,amssymb,amsthm,mathtools}
\usepackage{bm}

% -------------------------------------------------
% Layout
% -------------------------------------------------
\usepackage[a4paper,margin=2.5cm]{geometry}
\usepackage{hyperref}

% -------------------------------------------------
% Theorem-Umgebungen
% -------------------------------------------------
\newtheorem{definition}{Definition}[section]
\newtheorem{satz}[definition]{Satz}
\newtheorem{lemma}[definition]{Lemma}
\newtheorem{bemerkung}[definition]{Bemerkung}

% -------------------------------------------------
% Titel
% -------------------------------------------------
\title{Ein resonanter Operator auf Primzahlvierlingen\\
	\large Ein proto-spektales Modell im Sinne von Hilbert--Pólya}

\author{Thomas Hoffbauer\\
	Bamberg, Deutschland}

\date{\today}

% -------------------------------------------------
\begin{document}
	
	\maketitle
	
	\begin{abstract}
		Wir konstruieren einen nichtlokalen Operator auf der Menge von Primzahlvierlingen, der geometrische, arithmetische und phasenbasierte Strukturen kombiniert.
		Die resultierenden Eigenwertabstände zeigen eine nichttriviale Korrelation mit den Abständen der nichttrivialen Nullstellen der Riemannschen Zetafunktion.
		Das Modell verbindet diskrete Primzahlgeometrie mit kontinuierlicher spektraler Dynamik und liefert einen prototypischen Kandidaten im Kontext der Hilbert--Pólya-Idee.
	\end{abstract}
	
	\section{Einleitung}
	
	Die nichttrivialen Nullstellen der Riemannschen Zetafunktion werden seit langem mit Spektren selbstadjungierter Operatoren in Verbindung gebracht (Hilbert--Pólya-Vermutung).
	Parallel dazu zeigen Primzahlen in geeigneten Koordinaten (z.\,B.\ logarithmisch) eine überraschende Nähe zu spektralen Strukturen.
	
	In dieser Arbeit schlagen wir ein Modell vor, das Primzahlvierlinge als elementare Bausteine einer diskreten Geometrie interpretiert, auf der ein resonanter Operator definiert wird.
	
	\section{Primzahlvierlinge und Dehnungsfamilien}
	
	\begin{definition}[Primzahlvierling]
		Ein Primzahlvierling ist ein Tupel der Form
		\[
		Q(p) = (p,\, p+2,\, p+6,\, p+8),
		\]
		bestehend aus vier Primzahlen.
	\end{definition}
	
	\begin{definition}[Dehnungsfamilie]
		Wir definieren die Familie
		\[
		Q_k(p) = (p,\, p+2,\, p+6+12k,\, p+8+12k),
		\]
		die die mod-12-Struktur invariant erhält.
	\end{definition}
	
	Diese Struktur erlaubt eine stabile Klassifikation in die Restklassen
	\[
	E \equiv 1,\quad A \equiv 5,\quad B \equiv 7,\quad C \equiv 11 \pmod{12}.
	\]
	
	\section{Phasenstruktur}
	
	Wir interpretieren Primzahlen als Träger einer Phase:
	\begin{equation}
		\theta(p) = \alpha \log\log p + \theta_{\mathrm{EABC}}(p),
	\end{equation}
	wobei $\theta_{\mathrm{EABC}}$ eine diskrete Phase annimmt:
	\[
	\theta_{\mathrm{EABC}} \in \left\{0, \tfrac{\pi}{2}, \pi, \tfrac{3\pi}{2}\right\}.
	\]
	
	\begin{bemerkung}
		Die Überlagerung einer kontinuierlichen logarithmischen Phase mit einer diskreten mod-12-Struktur erzeugt eine intrinsisch interferierende Dynamik.
	\end{bemerkung}
	
	\section{Der Operator}
	
	Wir definieren einen Operator $H$ auf der Menge der Vierlinge durch
	
	\begin{equation}
		H_{ij} =
		\begin{cases}
			\log p_i, & i=j, \\[6pt]
			\dfrac{1}{|i-j|^\beta}
			\exp\bigl(i(\theta_i - \theta_j)\bigr), & i \neq j.
		\end{cases}
	\end{equation}
	
	\begin{bemerkung}
		Der Operator ist nichtlokal und enthält sowohl ein diagonales Potential als auch eine phasenmodulierte Kopplung.
		Für geeignete Wahl der Parameter kann $H$ als (nahezu) hermitesch interpretiert werden.
	\end{bemerkung}
	
	\section{Spektrale Eigenschaften}
	
	Seien $\lambda_n$ die Eigenwerte von $H$. Wir betrachten die normierten Abstände
	\[
	s_n = \frac{\lambda_{n+1} - \lambda_n}{\langle \lambda_{n+1} - \lambda_n \rangle}.
	\]
	
	\begin{satz}[Empirische Beobachtung]
		Die Verteilung der Abstände $s_n$ zeigt eine signifikante Ähnlichkeit mit der Wigner--Dyson-Verteilung (GUE-Typ).
	\end{satz}
	
	\section{Numerische Ergebnisse}
	
	Für moderate Systemgrößen zeigt sich:
	
	\begin{itemize}
		\item eine Korrelation der Abstände mit Zetazero-Abständen von etwa $0.2$--$0.25$,
		\item eine deutliche Abweichung von der Poisson-Statistik,
		\item eine Annäherung an GUE-ähnliches Verhalten.
	\end{itemize}
	
	\section{Interpretation}
	
	Das Modell legt nahe:
	
	\begin{itemize}
		\item Primzahlvierlinge bilden eine diskrete Geometrie,
		\item die Phase $\log\log p$ erzeugt spektrale Fluktuationen,
		\item nichtlokale Kopplung ist entscheidend für Quantenchaos.
	\end{itemize}
	
	\section{Diskussion und Ausblick}
	
	Das Modell ist kein Beweis der Riemannschen Vermutung, aber ein struktureller Kandidat im Sinne von Hilbert--Pólya.
	
	Offene Fragen:
	
	\begin{itemize}
		\item Exakte Wahl der Parameter $\alpha, \beta$,
		\item analytische Kontrolle des Spektrums,
		\item Zusammenhang mit bekannten Operatoren (Berry--Keating, Connes).
	\end{itemize}
	
	\section{Schlussbemerkung}
	
	Die hier vorgeschlagene Konstruktion verbindet arithmetische und spektrale Strukturen auf nichttriviale Weise und deutet darauf hin, dass Primzahlkonfigurationen als Träger eines tieferen dynamischen Systems interpretiert werden können.
	
	\bigskip
	
	\noindent
	\textbf{Schlüsselidee:}
	\emph{Primzahlen sind nicht nur Punkte, sondern Zustände eines resonanten Operators.}
	
\end{document}